\documentclass[11pt]{article}
\usepackage[margin=2.5cm]{geometry}
\usepackage{amsmath,amssymb,amsthm}
\usepackage{algorithm}
\usepackage{algpseudocode}
\usepackage{hyperref}

\newtheorem{proposition}{Proposition}
\newtheorem{remark}{Remark}
\newtheorem{lemma}{Lemma}

\DeclareMathOperator{\diag}{diag}
\DeclareMathOperator{\trace}{trace}
\DeclareMathOperator{\vMF}{vMF}
\DeclareMathOperator{\Ga}{Gamma}
\DeclareMathOperator*{\argmax}{arg\,max}

\newcommand{\X}{\mathbf{X}}
\newcommand{\U}{\mathbf{U}}
\newcommand{\V}{\mathbf{V}}
\newcommand{\EE}{\mathbf{E}}
\newcommand{\F}{\mathbf{F}}
\newcommand{\W}{\mathbf{W}}
\newcommand{\Q}{\mathbf{Q}}
\newcommand{\uu}{\mathbf{u}}
\newcommand{\vv}{\mathbf{v}}
\newcommand{\ww}{\mathbf{w}}
\newcommand{\f}{\mathbf{f}}
\newcommand{\z}{\mathbf{z}}
\newcommand{\I}{\mathbf{I}}
\newcommand{\one}{\mathbf{1}}
\newcommand{\Stiefel}[2]{\mathcal V_{#1,#2}}

\title{A Gibbs sampler for Optimal Multiple Co-Inertia Analysis}
\date{}

\begin{document}
\maketitle

\section{Probabilistic model}

We consider $K$ column-centered blocks $\X_1,\ldots,\X_K$, with $\X_k\in\mathcal M_{N,p_k}(\mathbb R)$, and a
latent space of dimension $H\leq\min_k p_k$. OMCOA minimizes
\begin{equation}
  \sum_{k=1}^K\bigl\|\X_k-\U\Sigma_k\V_k^\intercal\bigr\|^2
  \quad\text{subject to}\quad \U^\intercal\U=\I_H,\ \ \V_k^\intercal\V_k=\I_H,\ \ \Sigma_k=\diag\bigl(\sigma_k^{(1)},\ldots,\sigma_k^{(H)}\bigr).
  \label{eq:crit}
\end{equation}
A probabilistic counterpart of \eqref{eq:crit} is obtained by assuming a Gaussian noise:
\begin{equation}
  \X_k=\U\Sigma_k\V_k^\intercal+\EE_k,\qquad
  \EE_k\ \text{with i.i.d.\ entries}\ \mathcal N\bigl(0,\tau_k^{-1}\bigr),\qquad k=1,\ldots,K,
  \label{eq:model}
\end{equation}
where $\tau_k$ is the noise precision of block $k$. Let $\Stiefel{m}{H}=\{\W\in\mathcal M_{m,H}(\mathbb R):\W^\intercal\W=\I_H\}$
denote the Stiefel manifold. Since the blocks are centered, the latent components are constrained to be
centered as well: $\U\in\Stiefel{N}{H}$ with $\one_N^\intercal\U=0$.

\paragraph{Priors.} We use the following independent priors:
\begin{itemize}
  \item $\U$ and $\V_1,\ldots,\V_K$ are uniformly distributed on their Stiefel manifolds (non-informative);
  \item $\sigma_k^{(h)}\sim\mathcal N^+\bigl(0,\lambda^{-1}\bigr)$, a half-normal distribution with a small precision
        $\lambda$ (e.g.\ $\lambda=10^{-2}$); the positivity constraint removes the sign indeterminacy between
        $\sigma_k^{(h)}$ and $\vv_k^{(h)}$;
  \item $\tau_k\sim\Ga(a_0,b_0)$ with small hyperparameters (e.g.\ $a_0=b_0=10^{-3}$).
\end{itemize}

\section{Full conditional distributions}

\paragraph{Likelihood.} Because $\U$ and $\V_k$ have orthonormal columns and $\Sigma_k$ is diagonal,
\begin{equation}
  \bigl\|\X_k-\U\Sigma_k\V_k^\intercal\bigr\|^2
  =\|\X_k\|^2-2\trace\bigl(\Sigma_k\V_k^\intercal\X_k^\intercal\U\bigr)+\sum_{h=1}^H\bigl(\sigma_k^{(h)}\bigr)^2 .
  \label{eq:expand}
\end{equation}
Since the blocks are centered, each $\X_k$ carries $(N-1)p_k$ degrees of freedom, and the log-likelihood reads
\[
  \log p(\X_1,\ldots,\X_K\mid\cdot)=\sum_{k=1}^K\Bigl[\tfrac{(N-1)p_k}{2}\log\tau_k
  -\tfrac{\tau_k}{2}\bigl\|\X_k-\U\Sigma_k\V_k^\intercal\bigr\|^2\Bigr]+\text{const}.
\]
By \eqref{eq:expand}, $\U$ and $\V_k$ only appear in the linear term $\trace(\Sigma_k\V_k^\intercal\X_k^\intercal\U)$.
This is the key property: all the full conditionals belong to standard families.

\paragraph{Matrix von Mises--Fisher distribution.} A random matrix $\W\in\Stiefel{m}{H}$ follows the matrix
von Mises--Fisher distribution $\vMF(\F)$, with parameter $\F\in\mathcal M_{m,H}(\mathbb R)$, if its density with
respect to the uniform measure on $\Stiefel{m}{H}$ is proportional to $\exp\{\trace(\F^\intercal\W)\}$
\cite{hoff2009}. If $\F=\mathbf A\mathbf D\mathbf B^\intercal$ is the (thin) SVD of $\F$, the mode of $\vMF(\F)$ is the
polar factor $\mathbf A\mathbf B^\intercal$, and the distribution concentrates around it as the singular values of $\F$ increase.

\begin{proposition}[Full conditionals]\label{prop:cond}
Under model \eqref{eq:model} and the above priors:
\begin{enumerate}
  \item \textbf{Common space.} Let $\Q\in\mathcal M_{N,N-1}(\mathbb R)$ be an orthonormal basis of
        $\one_N^\perp$ and write $\U=\Q\tilde\U$ with $\tilde\U\in\Stiefel{N-1}{H}$. Then
        \[
          \tilde\U\mid\cdot\ \sim\ \vMF\bigl(\Q^\intercal\F_{\U}\bigr),\qquad
          \F_{\U}=\sum_{k=1}^K\tau_k\,\X_k\V_k\Sigma_k .
        \]
  \item \textbf{Block loadings.} For $k=1,\ldots,K$,
        \[
          \V_k\mid\cdot\ \sim\ \vMF\bigl(\F_{\V_k}\bigr),\qquad \F_{\V_k}=\tau_k\,\X_k^\intercal\U\Sigma_k .
        \]
  \item \textbf{Block weights.} For $k=1,\ldots,K$ and $h=1,\ldots,H$, let $d_k^{(h)}={\uu^{(h)}}^\intercal\X_k\vv_k^{(h)}$. Then
        \[
          \sigma_k^{(h)}\mid\cdot\ \sim\ \mathcal N^+\Bigl(\frac{\tau_k\,d_k^{(h)}}{\tau_k+\lambda},\ \frac{1}{\tau_k+\lambda}\Bigr),
        \]
        a normal distribution truncated to $(0,+\infty)$.
  \item \textbf{Noise precisions.} For $k=1,\ldots,K$,
        \[
          \tau_k\mid\cdot\ \sim\ \Ga\Bigl(a_0+\frac{(N-1)p_k}{2},\ b_0+\frac{1}{2}\bigl\|\X_k-\U\Sigma_k\V_k^\intercal\bigr\|^2\Bigr).
        \]
\end{enumerate}
\end{proposition}

\begin{proof}[Sketch of proof]
(1)--(2): by \eqref{eq:expand}, the conditional density of $\U$ is proportional to
\[
  \exp\Bigl\{\sum_k\tau_k\trace(\Sigma_k\V_k^\intercal\X_k^\intercal\U)\Bigr\}=\exp\{\trace(\F_{\U}^\intercal\U)\},
\]
and
$\trace(\F_{\U}^\intercal\Q\tilde\U)=\trace\bigl((\Q^\intercal\F_{\U})^\intercal\tilde\U\bigr)$; the same argument applies to $\V_k$
with $\trace(\Sigma_k\V_k^\intercal\X_k^\intercal\U)=\trace(\F_{\V_k}^\intercal\V_k)/\tau_k$.
(3): the terms of \eqref{eq:expand} involving $\sigma_k^{(h)}$ are $-\frac{\tau_k}{2}\bigl[(\sigma_k^{(h)})^2-2\sigma_k^{(h)}d_k^{(h)}\bigr]$;
adding the log-prior $-\frac{\lambda}{2}(\sigma_k^{(h)})^2$ and completing the square gives a Gaussian kernel, restricted to
$\sigma_k^{(h)}>0$.
(4): the Gamma prior is conjugate for the precision of a Gaussian likelihood.
\end{proof}

\subsection*{Detailed proof of item 2 (block loadings)}

Fix $k\in\{1,\ldots,K\}$ and denote by $\mu_k$ the uniform (Haar) probability measure on $\Stiefel{p_k}{H}$.

\paragraph{Step 1: only the likelihood of block $k$ matters.}
By Bayes' theorem, the full conditional density of $\V_k$ with respect to $\mu_k$ is proportional to the joint
density of the data and of all the parameters, seen as a function of $\V_k$ only. The priors are independent and the prior of $\V_k$ is uniform, i.e.\ its density with respect to $\mu_k$ is
constant. Moreover, given all the parameters, the blocks are independent and $\V_k$ only enters the distribution of
$\X_k$. Hence
\begin{equation}
  p(\V_k\mid\cdot)\ \propto\ p(\X_k\mid\U,\V_k,\Sigma_k,\tau_k)
  \ \propto\ \exp\Bigl\{-\frac{\tau_k}{2}\bigl\|\X_k-\U\Sigma_k\V_k^\intercal\bigr\|^2\Bigr\},
  \qquad \V_k\in\Stiefel{p_k}{H},
  \label{eq:step1}
\end{equation}
where the omitted factors do not depend on $\V_k$.

\paragraph{Step 2: expansion of the squared norm.}
Using $\|\mathbf A\|^2=\trace(\mathbf A^\intercal\mathbf A)$,
\[
  \bigl\|\X_k-\U\Sigma_k\V_k^\intercal\bigr\|^2
  =\|\X_k\|^2-2\trace\bigl(\X_k^\intercal\U\Sigma_k\V_k^\intercal\bigr)+\trace\bigl(\V_k\Sigma_k\U^\intercal\U\Sigma_k\V_k^\intercal\bigr).
\]
Since $\U^\intercal\U=\I_H$ and, by invariance of the trace under cyclic permutations, $\V_k^\intercal\V_k=\I_H$,
\[
  \trace\bigl(\V_k\Sigma_k\U^\intercal\U\Sigma_k\V_k^\intercal\bigr)=\trace\bigl(\Sigma_k^2\V_k^\intercal\V_k\bigr)=\trace(\Sigma_k^2)
  =\sum_{h=1}^H\bigl(\sigma_k^{(h)}\bigr)^2,
\]
which does not depend on $\V_k$: this is where the orthonormality constraint is essential. For the cross term, the
cyclic invariance and the invariance of the trace under transposition give
\[
  \trace\bigl(\X_k^\intercal\U\Sigma_k\V_k^\intercal\bigr)=\trace\bigl(\V_k^\intercal\X_k^\intercal\U\Sigma_k\bigr)
  =\trace\bigl((\X_k^\intercal\U\Sigma_k)^\intercal\V_k\bigr).
\]
Therefore
\begin{equation}
  -\frac{\tau_k}{2}\bigl\|\X_k-\U\Sigma_k\V_k^\intercal\bigr\|^2
  =\trace\bigl(\F_{\V_k}^\intercal\V_k\bigr)-\frac{\tau_k}{2}\Bigl(\|\X_k\|^2+\sum_{h=1}^H\bigl(\sigma_k^{(h)}\bigr)^2\Bigr),
  \qquad \F_{\V_k}=\tau_k\X_k^\intercal\U\Sigma_k .
  \label{eq:step2}
\end{equation}

\paragraph{Step 3: identification of the distribution.}
Plugging \eqref{eq:step2} into \eqref{eq:step1}, the second term of \eqref{eq:step2} is a constant and
\[
  p(\V_k\mid\cdot)\ \propto\ \exp\bigl\{\trace\bigl(\F_{\V_k}^\intercal\V_k\bigr)\bigr\},\qquad \V_k\in\Stiefel{p_k}{H}.
\]
The function $\V_k\mapsto\exp\{\trace(\F_{\V_k}^\intercal\V_k)\}$ is continuous on the compact set $\Stiefel{p_k}{H}$; it is therefore
bounded and integrable with respect to $\mu_k$, and the normalizing constant is finite and positive (it equals the
hypergeometric function ${}_0F_1\bigl(\tfrac{p_k}{2};\tfrac14\F_{\V_k}^\intercal\F_{\V_k}\bigr)$ \cite{hoff2009}). This is, by definition,
the density of the matrix von Mises--Fisher distribution $\vMF(\F_{\V_k})$ with respect to $\mu_k$, which proves item 2.
\hfill$\square$

\paragraph{Remarks on the proof.}
\begin{itemize}
  \item The diagonal structure of $\Sigma_k$ is not used: item 2 holds for any $H\times H$ matrix $\Sigma_k$. The orthonormality of
        $\U$ \emph{and} of $\V_k$ is used to make the quadratic term constant.
  \item If $\sigma_k^{(h)}=0$, the $h$-th column of $\F_{\V_k}$ vanishes and the conditional distribution does not favour any
        direction for $\vv_k^{(h)}$ (apart from its orthogonality to the other columns).
  \item Item 1 is proved in the same way: the log-likelihood depends on $\U$ through
        \[
          \sum_k\tau_k\trace\bigl((\X_k\V_k\Sigma_k)^\intercal\U\bigr)=\trace(\F_{\U}^\intercal\U),
        \]
        and the reparametrization $\U=\Q\tilde\U$, which maps $\Stiefel{N-1}{H}$ one-to-one onto the centered elements of
        $\Stiefel{N}{H}$ and preserves the uniform measure, gives
        $\trace(\F_{\U}^\intercal\Q\tilde\U)=\trace\bigl((\Q^\intercal\F_{\U})^\intercal\tilde\U\bigr)$.
\end{itemize}

\paragraph{Column-wise conditionals.}
Write $\W=\V_k$, $\F=\F_{\V_k}$ and $m=p_k$. Since $\trace(\F^\intercal\W)=\sum_{j=1}^H\f_j^\intercal\ww_j$, the conditional density of the
$j$-th column given the others is proportional to $\exp(\f_j^\intercal\ww_j)$. Given $\W_{-j}$, the column $\ww_j$ ranges over the unit
sphere of $\mathrm{span}(\W_{-j})^\perp$; writing $\ww_j=\mathbf N_j\z$ with $\mathbf N_j$ an orthonormal basis of this subspace and
$\z\in\mathbb S^{m-H}$, the uniform measure on $\Stiefel{m}{H}$ induces the uniform measure on $\mathbb S^{m-H}$ \cite{hoff2009}, and
$\f_j^\intercal\ww_j=(\mathbf N_j^\intercal\f_j)^\intercal\z$. Hence $\z\mid\W_{-j}\sim\vMF(\mathbf N_j^\intercal\f_j)$, which justifies the column-wise sweep
of Section~3.

\begin{lemma}[Mode of the matrix von Mises--Fisher distribution]\label{lem:mode}
Let $\F\in\mathcal M_{m,H}(\mathbb R)$ with thin SVD $\F=\mathbf A\mathbf D\mathbf B^\intercal$, $\mathbf D=\diag(\delta_1,\ldots,\delta_H)$, $\delta_j\geq0$.
For every $\W\in\Stiefel{m}{H}$, $\trace(\F^\intercal\W)\leq\sum_j\delta_j$, with equality for $\W=\mathbf A\mathbf B^\intercal$. If all $\delta_j>0$,
$\mathbf A\mathbf B^\intercal$ is the unique maximizer, hence the unique mode of $\vMF(\F)$.
\end{lemma}

\begin{proof}
Let $\mathbf C=\mathbf A^\intercal\W\mathbf B\in\mathcal M_{H,H}(\mathbb R)$. Then
$\trace(\F^\intercal\W)=\trace(\mathbf B\mathbf D\mathbf A^\intercal\W)=\trace(\mathbf D\mathbf A^\intercal\W\mathbf B)=\sum_j\delta_jC_{jj}$, and by the
Cauchy--Schwarz inequality $|C_{jj}|=|\mathbf a_j^\intercal\W\mathbf b_j|\leq\|\mathbf a_j\|\,\|\W\mathbf b_j\|=1$, since $\mathbf a_j$ and
$\mathbf b_j$ are unit vectors and $\W$ is an isometry. Hence $\trace(\F^\intercal\W)\leq\sum_j\delta_j$. For $\W=\mathbf A\mathbf B^\intercal$,
$\mathbf C=\I_H$ and equality holds. Conversely, if all $\delta_j>0$, equality forces $C_{jj}=1$, i.e.\ equality in Cauchy--Schwarz,
so $\W\mathbf b_j=\mathbf a_j$ for all $j$; since $\mathbf B$ is an orthogonal $H\times H$ matrix, $\W=\W\mathbf B\mathbf B^\intercal=\mathbf A\mathbf B^\intercal$.
\end{proof}

\noindent Applied to $\F_{\V_k}=\tau_k\X_k^\intercal\U\Sigma_k$, Lemma~\ref{lem:mode} shows that the mode of the conditional distribution of
$\V_k$ is the Procrustes solution of the SVD of $\X_k^\intercal\U\Sigma_k$ (the factor $\tau_k>0$ does not change the singular vectors),
which is exactly the OMCOA update of $\V_k$.

\section{Sampling from the matrix von Mises--Fisher distribution}

Direct sampling from $\vMF(\F)$ is not straightforward, but its column-wise conditionals are vector von
Mises--Fisher distributions \cite{hoff2009}. Let $\W_{-j}$ be $\W$ without its $j$-th column and
$\mathbf N_j\in\mathcal M_{m,m-H+1}(\mathbb R)$ an orthonormal basis of $\mathrm{span}(\W_{-j})^\perp$. Then
\[
  \ww_j=\mathbf N_j\z,\qquad \z\sim\vMF\bigl(\mathbf N_j^\intercal\f_j\bigr)\ \text{on the unit sphere } \mathbb S^{m-H},
\]
where $\f_j$ is the $j$-th column of $\F$. One sweep over the $H$ columns (in random order) is a valid Gibbs update
of $\W$, which is enough inside the Gibbs sampler. The vector distribution $\vMF(\f)$, with density
proportional to $\exp(\f^\intercal\z)$ on $\mathbb S^{m-1}$, is sampled with the rejection algorithm of Wood \cite{wood1994}
(Algorithm~\ref{alg:wood}).

\begin{algorithm}[ht]
\caption{Sampling $\z\sim\vMF(\f)$ on $\mathbb S^{m-1}$ (Wood, 1994)}\label{alg:wood}
\begin{algorithmic}[1]
  \State $\kappa\gets\|\f\|$,\quad $\boldsymbol\mu\gets\f/\kappa$
  \State $b\gets (m-1)\big/\bigl(2\kappa+\sqrt{4\kappa^2+(m-1)^2}\bigr)$,\quad $x_0\gets(1-b)/(1+b)$,\quad
         $c\gets\kappa x_0+(m-1)\log(1-x_0^2)$
  \Repeat
     \State draw $\zeta\sim\mathrm{Beta}\bigl(\tfrac{m-1}{2},\tfrac{m-1}{2}\bigr)$ and $\upsilon\sim\mathcal U(0,1)$
     \State $w\gets\bigl(1-(1+b)\zeta\bigr)\big/\bigl(1-(1-b)\zeta\bigr)$
  \Until{$\kappa w+(m-1)\log(1-x_0w)-c\geq\log\upsilon$}
  \State draw $\boldsymbol\xi\sim\mathcal N(\mathbf 0,\I_m)$, set $\boldsymbol\xi\gets\boldsymbol\xi-(\boldsymbol\mu^\intercal\boldsymbol\xi)\boldsymbol\mu$
         and $\boldsymbol\xi\gets\boldsymbol\xi/\|\boldsymbol\xi\|$
  \State \Return $\z=w\,\boldsymbol\mu+\sqrt{1-w^2}\,\boldsymbol\xi$
\end{algorithmic}
\end{algorithm}

\noindent The expression of $b$ used in Algorithm~\ref{alg:wood} is algebraically equivalent to Wood's original one,
$b=\bigl(-2\kappa+\sqrt{4\kappa^2+(m-1)^2}\bigr)/(m-1)$, but avoids a numerical cancellation when $\kappa$ is large,
which is the typical situation here since $\kappa$ scales with the noise precisions $\tau_k$.

\section{The Gibbs sampler}

\begin{algorithm}[ht]
\caption{Gibbs sampler for OMCOA}\label{alg:gibbs}
\begin{algorithmic}[1]
  \Require centered blocks $\X_1,\ldots,\X_K$; dimension $H$; number of iterations $T$; burn-in $T_0$;
           hyperparameters $a_0,b_0,\lambda$
  \State \textbf{Initialization:} $\U=\Q\tilde\U$, $\V_k$, $\Sigma_k$ from OMCOA (or MCOA, or at random);
         $\tau_k\gets(N-1)p_k\big/\|\X_k-\U\Sigma_k\V_k^\intercal\|^2$
  \For{$t=1,\ldots,T$}
     \State $\F_{\U}\gets\sum_{k}\tau_k\,\X_k\V_k\Sigma_k$;\quad update $\tilde\U$ by one column-wise sweep of $\vMF(\Q^\intercal\F_{\U})$;\quad $\U\gets\Q\tilde\U$
     \For{$k=1,\ldots,K$}
        \State $\F_{\V_k}\gets\tau_k\,\X_k^\intercal\U\Sigma_k$;\quad update $\V_k$ by one column-wise sweep of $\vMF(\F_{\V_k})$
        \For{$h=1,\ldots,H$}
           \State $d_k^{(h)}\gets{\uu^{(h)}}^\intercal\X_k\vv_k^{(h)}$;\quad draw
                  $\sigma_k^{(h)}\sim\mathcal N^+\bigl(\tau_kd_k^{(h)}/(\tau_k+\lambda),\,1/(\tau_k+\lambda)\bigr)$
        \EndFor
        \State draw $\tau_k\sim\Ga\bigl(a_0+(N-1)p_k/2,\ b_0+\|\X_k-\U\Sigma_k\V_k^\intercal\|^2/2\bigr)$
     \EndFor
     \If{$t>T_0$} store $\bigl(\U,\V_1,\ldots,\V_K,\Sigma_1,\ldots,\Sigma_K,\tau_1,\ldots,\tau_K\bigr)$ \EndIf
  \EndFor
  \State \textbf{Post-processing:} align each stored draw on a reference solution (Section~\ref{sec:align})
  \State \Return the aligned draws
\end{algorithmic}
\end{algorithm}

\section{Post-processing: label switching and signs}\label{sec:align}

The posterior distribution is invariant under any permutation $\pi$ of the $H$ dimensions applied jointly to
$\U$, $\V_1,\ldots,\V_K$ and $\Sigma_1,\ldots,\Sigma_K$, and under a joint change of sign of $\uu^{(h)}$ and
$\vv_1^{(h)},\ldots,\vv_K^{(h)}$. Posterior summaries are therefore computed after aligning each draw on a reference
$\U_{\mathrm{ref}}$ (in practice the OMCOA solution). For each draw, with $\mathbf C=\U_{\mathrm{ref}}^\intercal\U$,
\[
  \hat\pi=\argmax_{\pi\in\mathfrak S_H}\sum_{h=1}^H\bigl|C_{h,\pi(h)}\bigr|,\qquad
  s_h=\operatorname{sign}\bigl(C_{h,\hat\pi(h)}\bigr),
\]
and the columns of $\U$ and $\V_k$ are reordered by $\hat\pi$ and multiplied by $s_h$, while the diagonal of
$\Sigma_k$ is reordered by $\hat\pi$ (it stays positive). Posterior means and credible intervals of $\sigma_k^{(h)}$
are then obtained from the empirical distribution of the aligned draws.

\section{Link with OMCOA}

\begin{remark}[OMCOA as iterated conditional modes]
The modes of the full conditionals of Proposition~\ref{prop:cond} are
\[
  \U=\mathbf A\mathbf B^\intercal\ \text{with}\ \F_{\U}=\mathbf A\mathbf D\mathbf B^\intercal,\qquad
  \V_k=\mathbf A_k\mathbf B_k^\intercal\ \text{with}\ \F_{\V_k}=\mathbf A_k\mathbf D_k\mathbf B_k^\intercal,\qquad
  \sigma_k^{(h)}=\frac{\tau_k}{\tau_k+\lambda}\,d_k^{(h)} .
\]
When $\lambda\to0$ and the precisions are equal ($\tau_1=\cdots=\tau_K$), these are exactly the OMCOA updates:
Procrustes rotation of $\X\mathbf H^\intercal$ for $\U$ (with $\mathbf H=[\Sigma_1\V_1^\intercal|\cdots|\Sigma_K\V_K^\intercal]$),
Procrustes rotation of $\X_k^\intercal\U\Sigma_k$ for $\V_k$, and $\Sigma_k=\diag(\U^\intercal\X_k\V_k)$. OMCOA is therefore the
iterated conditional modes (ICM) algorithm \cite{besag1986} associated with this posterior, and the Gibbs sampler is its
stochastic counterpart.
\end{remark}

\begin{remark}[Block weighting]
With block-specific precisions, $\F_{\U}=\sum_k\tau_k\X_k\V_k\Sigma_k$ weights each block by its estimated noise
precision, so that noisier blocks contribute less to the common space. To reproduce exactly the OMCOA criterion
\eqref{eq:crit}, a common precision $\tau_1=\cdots=\tau_K=\tau$ can be used, with
$\tau\mid\cdot\sim\Ga\bigl(a_0+\sum_k(N-1)p_k/2,\ b_0+\sum_k\|\X_k-\U\Sigma_k\V_k^\intercal\|^2/2\bigr)$.
\end{remark}

\begin{remark}[What the sampler adds]
Besides point estimates close to the OMCOA solution, the sampler provides the posterior distribution of the block
weights $\sigma_k^{(h)}$ and of the noise levels $\tau_k^{-1}$. Credible intervals on $\sigma_k^{(h)}$ can be used to assess
whether a dimension is common to all blocks, shared by some of them, or specific to one block.
\end{remark}

\begin{remark}[Practical aspects]
Starting from the OMCOA solution shortens the burn-in. The column-wise sweeps make successive draws of $\U$ and
$\V_k$ correlated, so convergence should be checked on trace plots of the $\sigma_k^{(h)}$. Because the posterior is
strongly concentrated when the noise is small, the chain rarely jumps between the $H!$ equivalent modes; the alignment step
remains necessary nonetheless, in particular when several dimensions have close weights.
\end{remark}

\begin{thebibliography}{9}
\bibitem{besag1986} J.~Besag (1986). On the statistical analysis of dirty pictures.
  \emph{Journal of the Royal Statistical Society, Series B}, 48(3), 259--302.
\bibitem{hoff2007} P.~D.~Hoff (2007). Model averaging and dimension selection for the singular value decomposition.
  \emph{Journal of the American Statistical Association}, 102(478), 674--685.
\bibitem{hoff2009} P.~D.~Hoff (2009). Simulation of the matrix Bingham--von Mises--Fisher distribution, with
  applications to multivariate and relational data. \emph{Journal of Computational and Graphical Statistics}, 18(2), 438--456.
\bibitem{wood1994} A.~T.~A.~Wood (1994). Simulation of the von Mises Fisher distribution.
  \emph{Communications in Statistics -- Simulation and Computation}, 23(1), 157--164.
\end{thebibliography}

\end{document}
